\documentclass[11pt]{article}
% ===========================================================================
%  Shared preamble for the AP Calculus AB exam-preparation worksheets.
%  Input right after \documentclass. Files beginning with "_" are not compiled
%  on their own. Navy and gold, to match the course's lesson pages.
% ===========================================================================
\usepackage[letterpaper,top=2.3cm,bottom=2.2cm,left=1.7cm,right=1.7cm,headheight=28pt,footskip=30pt]{geometry}
\usepackage{amsmath,amssymb}
\usepackage[T1]{fontenc}
\usepackage[utf8]{inputenc}
\usepackage{lmodern}
\usepackage{xcolor}
\usepackage[most]{tcolorbox}
\usepackage{enumitem}
\usepackage{fancyhdr}
\usepackage{array}
\usepackage{tikz}
\usepackage{pgfplots}
\pgfplotsset{compat=1.18}

\definecolor{navy}{HTML}{161E54}
\definecolor{ink}{HTML}{27357E}
\definecolor{gold}{HTML}{B98A00}
\definecolor{goldbg}{HTML}{FFFDF3}
\definecolor{inkbg}{HTML}{F6F7FD}
\definecolor{rule}{HTML}{C9CFF0}
\definecolor{okgreen}{HTML}{1B7A44}
\definecolor{okbg}{HTML}{F1F8F3}

\pagestyle{fancy}
\fancyhf{}
\fancyhead[L]{\footnotesize Name: \underline{\hspace{4cm}}}
\fancyhead[C]{\textbf{\color{navy}Integration Academy}\\[-2pt]{\scriptsize Calculus AB \textperiodcentered{} AP\textsuperscript{\textregistered} Exam Preparation}}
\fancyfoot[L]{\tiny\color{black!55}\parbox[t]{0.78\textwidth}{AP\textsuperscript{\textregistered} is a registered trademark of the College Board, which was not involved in the production of, and does not endorse, this material. \textcopyright{} Integration Academy.}}
\fancyfoot[R]{\scriptsize\thepage}
\renewcommand{\headrulewidth}{1.2pt}
\renewcommand{\headrule}{\hbox to\headwidth{\color{gold}\leaders\hrule height \headrulewidth\hfill}}
\renewcommand{\footrulewidth}{0pt}
\setlength{\parindent}{0pt}
\setlength{\parskip}{4pt}

% Title block: \wstitle{code}{title}{unit line}
\newcommand{\wstitle}[3]{\begin{center}{\small\bfseries\color{gold}WORKSHEET #1}\\[3pt]{\LARGE\bfseries\color{navy}#2}\\[3pt]{\small\color{black!60}#3}\end{center}\vspace{2pt}}
% Section heading
\newcommand{\sech}[1]{\par\vspace{10pt}{\large\bfseries\color{navy}#1}\par\nobreak\vspace{2pt}{\color{rule}\hrule height0.8pt}\vspace{7pt}}
% A level tag printed beside a group of questions
\newcommand{\level}[1]{\par\vspace{6pt}{\footnotesize\bfseries\color{gold}\MakeUppercase{#1}}\par\nobreak\vspace{2pt}}

\newtcolorbox{ideabox}{colframe=ink,colback=inkbg,coltitle=white,colbacktitle=ink,fonttitle=\bfseries\small,title={Key ideas},arc=2pt,boxrule=0.8pt,left=8pt,right=8pt}
\newtcolorbox{example}[1]{colframe=rule,colback=white,coltitle=navy,colbacktitle=inkbg,fonttitle=\bfseries\small,title={#1},arc=2pt,boxrule=0.8pt,borderline west={3pt}{0pt}{ink},left=10pt,right=8pt}
\newtcolorbox{frbox}[1]{colframe=gold,colback=goldbg,coltitle=white,colbacktitle=gold,fonttitle=\bfseries\small,title={#1},arc=2pt,boxrule=0.8pt,breakable,left=8pt,right=8pt}
\newtcolorbox{solbox}{colframe=okgreen,colback=okbg,coltitle=white,colbacktitle=okgreen,fonttitle=\bfseries\small,title={Worked solutions},arc=2pt,boxrule=0.8pt,breakable,left=8pt,right=8pt}
\newtcolorbox{scorebox}[1]{colframe=navy,colback=white,coltitle=white,colbacktitle=navy,fonttitle=\bfseries\small,title={#1},arc=2pt,boxrule=0.8pt,breakable,left=8pt,right=8pt}

\newcommand{\soln}{\par\smallskip\textbf{\color{ink}Solution.}\ }
\newcommand{\result}[1]{\par\smallskip\textbf{#1}}
% blank working space of a given height
\newcommand{\work}[1]{\par\vspace{#1}}
% a point value in the margin of a scoring guide
\newcommand{\pts}[1]{\hfill{\footnotesize\bfseries\color{gold}[#1]}}

% Coordinate grid for graph questions: \begin{apgraph}{xmin}{xmax}{ymin}{ymax} ... \end{apgraph}
\newenvironment{apgraph}[4]{\begin{center}\begin{tikzpicture}\begin{axis}[
  axis lines=middle, xmin=#1, xmax=#2, ymin=#3, ymax=#4,
  grid=both, grid style={black!12}, xtick distance=1, ytick distance=1,
  ticklabel style={font=\scriptsize}, xlabel={$x$}, ylabel={$y$},
  width=10cm, height=7.2cm, samples=80, clip=false,
  every axis plot/.append style={line width=1.1pt, color=ink}]}{\end{axis}\end{tikzpicture}\end{center}}
\newcommand{\closeddot}[2]{\addplot[only marks, mark=*, mark size=2.2pt, color=ink] coordinates {(#1,#2)};}
\newcommand{\opendot}[2]{\addplot[only marks, mark=*, mark size=2.2pt, color=ink, mark options={fill=white}] coordinates {(#1,#2)};}

% Multiple-choice options: four across, or stacked when the options are long.
\newcommand{\opts}[4]{\par\vspace{3pt}\begin{tabular}{@{}p{0.235\linewidth}p{0.235\linewidth}p{0.235\linewidth}p{0.235\linewidth}@{}}(A)\ #1 & (B)\ #2 & (C)\ #3 & (D)\ #4\end{tabular}\par\vspace{9pt}}
\newcommand{\optsv}[4]{\par\vspace{3pt}\begin{tabular}{@{}l@{\ \ }p{0.9\linewidth}@{}}(A) & #1\\(B) & #2\\(C) & #3\\(D) & #4\end{tabular}\par\vspace{9pt}}

\fancyhead[R]{\footnotesize\color{black!60}Worksheet 3.1}
\begin{document}

\wstitle{3.1}{Functions Inside Functions: The Chain Rule}{Unit 3: Composite, Implicit and Inverse Functions}

\begin{ideabox}
\begin{itemize}[leftmargin=*,itemsep=2pt,topsep=2pt]
  \item \textbf{Chain rule}: $\dfrac{d}{dx}\Big[f\big(g(x)\big)\Big]=f'\big(g(x)\big)\cdot g'(x)$: the derivative of the outside, evaluated at the inside, times the derivative of the inside.
  \item Work from the outside in. With three layers there are three factors.
  \item $\sin^3x$ means $(\sin x)^3$: the outside function is the cube.
  \item With a table, $f'$ is evaluated at $g(a)$, not at $a$.
\end{itemize}
\end{ideabox}

\sech{Worked examples}
\begin{example}{Example 1 \textperiodcentered{} A power of a polynomial}
Find $\dfrac{d}{dx}\left[\left(2x^3+1\right)^5\right]$.\soln $5\left(2x^3+1\right)^4\cdot6x^2$.\result{$30x^2\left(2x^3+1\right)^4$.}
\end{example}

\begin{example}{Example 2 \textperiodcentered{} A root of a trigonometric function}
Find $\dfrac{d}{dx}\left[\sqrt{3+\cos x}\right]$.\soln $\dfrac12(3+\cos x)^{-1/2}\cdot(-\sin x)$.\result{$-\dfrac{\sin x}{2\sqrt{3+\cos x}}$.}
\end{example}

\begin{example}{Example 3 \textperiodcentered{} From values}
Let $h(x)=f\big(g(x)\big)$, with $g(3)=2$, $g'(3)=-4$ and $f'(2)=5$. Find $h'(3)$.\soln $h'(3)=f'\big(g(3)\big)\,g'(3)=f'(2)\cdot(-4)$.\result{$h'(3)=-20$.}
\end{example}

\newpage
\sech{Practice}

\level{Core}
\begin{enumerate}[leftmargin=*,itemsep=2pt]
  \item Find $\dfrac{d}{dx}[\sin(5x)]$.\work{2.3cm}
  \item Find $\dfrac{d}{dx}\left[e^{x^2}\right]$.\work{2.3cm}
  \item Find $\dfrac{d}{dx}\left[\left(x^2+3x\right)^4\right]$.\work{2.3cm}
  \item Find $\dfrac{d}{dx}[\ln(4x+1)]$.\work{2.3cm}
\end{enumerate}

\level{Exam level}
\begin{enumerate}[leftmargin=*,itemsep=2pt,start=5]
  \item Find $\dfrac{d}{dx}\left[\cos^3x\right]$.\work{2.8cm}
  \item Find $\dfrac{d}{dx}\left[e^{-2x}\sin(3x)\right]$.\work{2.8cm}
  \item Find $\dfrac{d}{dx}\left[\tan\!\left(\sqrt{x}\right)\right]$.\work{2.8cm}
  \item Let $f(x)=\sqrt{g(x)}$, where $g(4)=25$ and $g'(4)=-10$. Find $f'(4)$.\work{2.8cm}
\end{enumerate}

\begin{frbox}{Question 9 \textperiodcentered{} Free-response practice}
Let $f(x)=\left(x^2-4\right)^3$.
\begin{enumerate}[label=(\alph*),leftmargin=*,itemsep=2pt]
  \item Find $f'(x)$.\work{2.6cm}
  \item Find all values of $x$ at which the graph of $f$ has a horizontal tangent.\work{2.6cm}
  \item Write an equation of the line tangent to the graph of $f$ at $x=1$.\work{2.6cm}
\end{enumerate}
\end{frbox}

\newpage
\begin{solbox}
\begin{enumerate}[leftmargin=*,itemsep=4pt]
  \item $5\cos(5x)$.
  \item $2xe^{x^2}$.
  \item $4\left(x^2+3x\right)^3(2x+3)$.
  \item $\dfrac{4}{4x+1}$.
  \item $3\cos^2x\cdot(-\sin x)=-3\cos^2x\sin x$.
  \item $-2e^{-2x}\sin(3x)+3e^{-2x}\cos(3x)=e^{-2x}\big(3\cos(3x)-2\sin(3x)\big)$.
  \item $\sec^2\!\left(\sqrt x\right)\cdot\dfrac{1}{2\sqrt x}$.
  \item $f'(4)=\dfrac{g'(4)}{2\sqrt{g(4)}}=\dfrac{-10}{10}=-1$.
  \item \begin{enumerate}[label=(\alph*),leftmargin=*,itemsep=3pt]
    \item $3\left(x^2-4\right)^2\cdot2x=6x\left(x^2-4\right)^2$.
    \item $x=0$, $x=2$ and $x=-2$.
    \item $f(1)=-27$ and $f'(1)=6(9)=54$: $y+27=54(x-1)$.
  \end{enumerate}
\end{enumerate}
\end{solbox}

\end{document}
